
\documentclass[UTF8,12pt,a4paper]{ctexart}
\usepackage[margin=2.2cm]{geometry}
\usepackage{amsmath,amssymb,mathtools}
\usepackage{fontspec}
\usepackage{xcolor}
\usepackage{enumitem}
\usepackage{hyperref}
\usepackage{fancyvrb}
\usepackage{fvextra}
\usepackage{microtype}
\usepackage{titlesec}
\usepackage{fancyhdr}
\usepackage{setspace}

\hypersetup{
  colorlinks=true,
  linkcolor=blue!55!black,
  urlcolor=blue!55!black
}

\setstretch{1.18}
\setlist[itemize]{leftmargin=1.8em,itemsep=0.25em,topsep=0.3em}
\setlist[enumerate]{leftmargin=2.1em,itemsep=0.25em,topsep=0.3em}
\titleformat{\section}{\Large\bfseries}{\thesection}{0.8em}{}
\titleformat{\subsection}{\large\bfseries}{\thesubsection}{0.8em}{}
\titleformat{\subsubsection}{\normalsize\bfseries}{\thesubsubsection}{0.8em}{}
\pagestyle{fancy}
\fancyhf{}
\fancyfoot[C]{\thepage}
\renewcommand{\headrulewidth}{0pt}

\DefineVerbatimEnvironment{GuideVerbatim}{Verbatim}{
  breaklines=true,
  breakanywhere=true,
  fontsize=\small,
  frame=single,
  framesep=3mm,
  rulecolor=\color{gray!35},
  bgcolor=gray!7
}

\newcommand{\GuideInlineCode}[1]{\texttt{\detokenize{#1}}}

\begin{document}

\section*{路线 4：QAOA / VQA 标准量子门模型}

QAOA 是量子门模型里最常见的组合优化算法路线。它的核心不是“量子电脑自动给答案”，而是构造一个带参数的量子线路，让低能 bitstring 出现概率更高，再用经典优化器调这些参数。

第一遍学 QAOA，要抓住三个东西：cost Hamiltonian、mixer、classical optimizer。

\subsection*{QAOA 在做什么}

QUBO 给每个 bitstring 一个能量。

QAOA 希望准备一个量子态，使得测量时更容易测到低能 bitstring。

它不是直接输出一个确定解，而是输出一个概率分布：

\begin{GuideVerbatim}
00110: 0.21
10100: 0.17
00111: 0.09
...
\end{GuideVerbatim}

然后我们从测量结果里解码、检查约束、挑 best feasible solution。

\subsection*{标准 QAOA 结构}

QAOA 的状态一般写成：

\[
\begin{aligned}
|\psi(gamma, beta)\rangle\\
= \prod_{l} \exp(-i beta_{l} H_{M}) \exp(-i gamma_{l} H_{C}) |+\rangle^{\otimes n}
\end{aligned}
\]

你不用被公式吓到。它只有三层意思：

\texttt{|+>\textasciicircum{}n}

一开始让所有 bitstring 大致均匀叠加。可以理解成“所有候选解先都在场”。

\texttt{H\_C}

Cost Hamiltonian。它来自 QUBO/Ising，负责给不同 bitstring 加上和成本有关的相位。

\texttt{H\_M}

Mixer Hamiltonian。它负责在 bitstring 之间混合，让概率流动起来。

\texttt{gamma} 和 \texttt{beta}

可调参数。经典优化器会不断试参数，观察期望能量，寻找更好的概率分布。

\subsection*{Cost Hamiltonian}

如果 QUBO 是：

\[
E(x) = c + \sum_{i} q_{i} x_{i} + \sum_{ij} q_{ij} x_{i} x_{j}
\]

转成 Ising 后，就可以构造 cost Hamiltonian：

\[
H_{C} = C + \sum_{i} h_{i} Z_{i} + \sum_{ij} J_{ij} Z_{i} Z_{j}
\]

\texttt{Z\_i} 可以理解成读取第 \texttt{i} 个 qubit 的 0/1 信息。某个 bitstring 的能量越低，我们越希望它最后被测出来的概率越高。

\subsection*{Mixer 是什么}

标准 QAOA 常用 X mixer：

\[
H_{M} = \sum_{i} X_{i}
\]

它的直觉是：每个 bit 都可以翻转。所以标准 X mixer 会在整个 \texttt{2\textasciicircum{}n} 搜索空间里移动。

这很通用，但也有问题：它会移动到非法状态。例如 exactly-one 问题里，从 \texttt{100} 翻一个 bit 可能变成 \texttt{110} 或 \texttt{000}，都不合法。

所以标准 QAOA 常常还需要 penalty 来处理约束。路线 5 的 constrained mixer，就是在改 mixer。

\subsection*{VQA 是更大的框架}

VQA 是 Variational Quantum Algorithm，QAOA 是 VQA 的一个特化版本。

VQA 的共同流程：

\begin{GuideVerbatim}
1. 设计参数化量子线路 ansatz
2. 在量子后端运行线路
3. 测量得到样本或期望值
4. 经典优化器根据结果更新参数
5. 重复直到预算耗尽或收敛
\end{GuideVerbatim}

近端量子算法的典型样子就是这种量子-经典闭环。

\subsection*{p 层是什么意思}

QAOA 的 \texttt{p} 是交替层数：

\[
\begin{aligned}
p = 1: cost \to mixer\\
p = 2: cost \to mixer \to cost \to mixer\\
...
\end{aligned}
\]

\texttt{p} 越大，表达能力通常越强，但线路越深、参数越多、优化更难、噪声更重。

比赛里不要默认“p 越大越好”。小规模 demo 可以比较：

\[
p = 1, 2, 3
\]

记录 best feasible objective、feasible ratio、optimizer evaluations、runtime。

\subsection*{Classical optimizer 的角色}

QAOA 的参数不是量子设备自己学出来的。经典优化器负责试参数。

常见优化器：

\begin{itemize}
\item COBYLA
\item Nelder-Mead
\item SPSA
\item gradient-based methods
\item multi-start random search
\end{itemize}

QAOA 的难点之一就是参数地形可能很不友好：有局部最优、shot noise、barren plateau、参数周期性。

所以参数初始化也很重要。随机初始化是 baseline；更好的方法包括 linear ramp、warm-start parameters、复用小规模问题参数。

\subsection*{Measurement 和后处理}

QAOA 最后测量会得到 counts：

\begin{GuideVerbatim}
00101: 120
00111: 84
10101: 31
...
\end{GuideVerbatim}

后处理步骤和退火很像：

\begin{GuideVerbatim}
counts
  -> bitstrings
  -> QUBO energy
  -> decode
  -> feasibility check
  -> original objective
  -> best feasible
\end{GuideVerbatim}

不要只报告期望能量。比赛里最终要的是可解释的业务解。

\subsection*{QAOA 适合什么}

小规模组合优化 demo。

QAOA 的 qubit 数和线路深度限制很明显，所以它适合拿小实例做算法展示。

和 classical baseline 对比。

比如 exact small case、random、SA、standard QAOA、constrained QAOA。

展示量子结构。

如果你能解释 cost layer、mixer、测量分布，而不是只说“调用 QAOA”，评委会更容易相信你懂模型。

\subsection*{比赛时能改哪些地方}

Cost Hamiltonian

QUBO 建得好不好，会直接影响 QAOA。

Mixer

标准 X mixer 是 baseline。问题有 exactly-one 或 fixed-k 结构时，可以改成 XY mixer 或 problem-specific mixer。

Initial state

从均匀叠加开始不是唯一选择。可以用 feasible initial state 或 warm-start state。

参数初始化

random、linear ramp、warm-start、multi-start 都可以比较。

Optimizer

不同 optimizer 对 shot noise 和局部最优的表现不同。

后处理

QAOA 输出的是概率分布。repair、可行性优先排序、local improvement 都有空间。

\subsection*{常见坑}

把 QAOA 当成大规模万能 solver。

当前近端环境下，QAOA 更适合小规模算法展示和结构对比。

忽略 constraints。

标准 X mixer 会走到非法状态。约束如果只靠 penalty，penalty 质量仍然决定可行性。

只看期望能量。

期望能量低不代表测量样本里有高质量可行解。要看 best feasible 和 feasible ratio。

参数优化预算不公平。

和 SA 比较时，要说明 shots、optimizer evaluations、runtime。否则对比不公平。

线路深度过深。

真实硬件或 noisy simulator 上，深线路可能被噪声毁掉。

\subsection*{和其他路线的关系}

路线 1 把原问题变成 cost Hamiltonian。

路线 2 决定 penalty 形式和可行性检查。

路线 3 可以作为 QAOA 的经典/退火 baseline。

路线 5 是 QAOA 的结构升级：改 mixer、改 initial state、改 feasible subspace。

路线 6 可以先把大问题拆成小 QUBO，再用 QAOA 求子问题。

\subsection*{共读任务}

用一句话回答：

\begin{GuideVerbatim}
QAOA 的 cost layer 在奖励什么？mixer layer 又在允许什么？
\end{GuideVerbatim}

如果能答出“cost layer 来自目标能量，mixer layer 决定搜索空间里的移动方式”，你就抓住了 QAOA 的骨架。

\subsection*{延伸阅读}

\begin{itemize}
\item \texttt{literature/requirements/04\_qaoa\_vqa\_requirements.md}
\item Farhi, Goldstone, Gutmann, QAOA original paper
\item Zhou et al., QAOA performance and parameter mechanisms
\item Cerezo et al., VQA review
\end{itemize}

\end{document}
